\section{Discussion}
The results should be interpreted as one empirical comparison on one financial
index over one historical period. Forecasting accuracy and trend recovery are
different goals, and a method that forecasts well over the selected test window
need not recover a latent economic trend. Validation-based selection can also
overfit the validation period, especially when the candidate grid is wide or
the market regime changes between validation and test.

Guerrero smoothness is a model-side index defined only for finite-difference
penalized least-squares smoothers. It maps a smoothing penalty \(\lambda\) into
a comparable scale \(S_d(\lambda; N)\) once the sample size \(N\) and
difference order \(d\) are fixed. By contrast, realized roughness is an output
diagnostic \(R_d(\hat{\tau}) = \|D^d\hat{\tau}\|_2^2\) that can be computed for
any fitted trend. Therefore, moving-average, exponential-smoothing, and
polynomial benchmarks do not have a native Guerrero smoothness index, although
their fitted trends can still be compared through realized roughness.

Future sensitivity checks should vary the time period, include additional
tickers, and compare against synthetic experiments where the true trend is
known. Additional baselines such as Kalman filters, structural time-series
models, ARIMA models, STL/LOESS, smoothing splines, and structural-break models
would also make the comparison broader
\citep{kalman1960new,harvey1989structural,box2008time,cleveland1979lowess,
cleveland1990stl,wahba1990splines}.
